% ---------------------------------------------------------------------------
% Section: the distance-trimmed time-unrolled walk encoding.
% Paste-ready. TODO markers point to labels of the main text.
% ---------------------------------------------------------------------------

\section{A distance-trimmed time-unrolled encoding of grid walks}
\label{sec:trimmed-steps}

Let $G=(V,E)$ be the four-neighbour grid graph on a $g\times g$ board,
$n=g^2=|V|$, with source $s$ (top-left) and target $t$ (bottom-right).
Write $d(u,v)$ for the graph (Manhattan) distance. Recall that $G$ is
bipartite under the checkerboard two-colouring: every edge joins cells of
different colours, hence a walk of length $\ell$ from $s$ can only end on
a cell $v$ whose distance $d(s,v)$ has the same parity as $\ell$
($\ell \equiv d(s,v) \pmod 2$; e.g.\ a cell at even distance is reachable
only in an even number of moves).

\paragraph{Variables and transitions.}
Fix the horizon $L = n-1$, the maximum number of edges of any simple
path. Introduce multi-valued variables $X_0,\dots,X_L$ with intended
reading ``$X_i$ = the cell occupied at time $i$''. The \emph{absorbing
transition relation} is
\[
  T(u,v) \iff \bigl(u\neq t \wedge \{u,v\}\in E\bigr) \;\lor\; u=v=t ,
\]
i.e.\ every step is a real move, except that the target, once reached,
repeats forever (and can never be left). (There is no general waiting move; this is what makes
the parity argument below available.) The formula is the chain
\[
  \Phi \;=\; \bigwedge_{i=0}^{L-1} T(X_i, X_{i+1}),
\]
endpoint pinning being absorbed into the domains below.

\paragraph{From the relation to signed clauses.}
$T$ is not a primitive of the constraint language; it is unrolled value
by value. Writing $[S]{:}X$ for the signed literal ``$X\in S$'', each
conjunct $T(X_i,X_{i+1})$ becomes one clause per value $u\in D_i$,
\[
  [\,D_i\setminus\{u\}\,]{:}X_i \;\vee\;
  [\,\mathrm{succ}(u)\cap D_{i+1}\,]{:}X_{i+1},
  \qquad
  \mathrm{succ}(u)=\begin{cases}
    N(u) & u\neq t,\\
    \{t\} & u=t,
  \end{cases}
\]
i.e.\ the implication $X_i{=}u \Rightarrow X_{i+1}\in\mathrm{succ}(u)$
written as a disjunction. The edge set $E$ never appears as an object:
adjacency ($N(u)$ = the up/down/left/right cells of $u$) is materialized
inside the enumerated successor lists. A guard $[\emptyset]{:}X_i$
(empty when $D_i=\{u\}$ is a singleton) is dropped, leaving a unary
clause. For the trimmed $2\times 2$ instance the whole formula is four
clauses:
\[
  [1,2]{:}X_1, \qquad
  [2]{:}X_1 \vee [3]{:}X_2, \qquad
  [1]{:}X_1 \vee [3]{:}X_2, \qquad
  [3]{:}X_3,
\]
(step $0$: $D_0=\{0\}$ forces $X_1\in N(0)\cap D_1=\{1,2\}$; step $1$:
from $u=1$ or $u=2$ the only admissible successor in $D_2=\{3\}$ is the
target; step $2$: absorption). Optionally,
\textsc{AllDifferent} clauses over the non-target positions restrict the
models to simple paths; we omit them (\texttt{allow\_cycles}) because
they connect every pair of step variables and destroy the chain
structure — simplicity is instead recovered downstream (loop erasure) or
left to the learner.

\paragraph{Domain trimming.}
The untrimmed encoding gives every $X_i$ the full domain $V$. The
trimmed encoding declares instead
\[
  D_i \;=\;
  \bigl\{\, v \in V\setminus\{t\} \;:\;
      \underbrace{d(s,v)\le i}_{\text{forward}},\;
      \underbrace{i \equiv d(s,v) \!\!\pmod 2}_{\text{parity}},\;
      \underbrace{d(v,t)\le L-i}_{\text{backward}} \,\bigr\}
  \;\cup\;
  \bigl\{\, t \;:\; i \ge d(s,t) \,\bigr\}.
\]
In particular $D_0=\{s\}$ and $D_L=\{t\}$: the endpoint constraints are
domain facts, not clauses.

\paragraph{Reading the three tests.}
Since $D_L=\{t\}$, the walk must have \emph{arrived by the deadline}
$L$; at time $i$ it therefore has exactly $L-i$ transitions left. The
three tests bound where it can possibly stand at that moment:
\begin{itemize}
\item \emph{Forward} ($d(s,v)\le i$): the past constrains the present —
  in $i$ moves from $s$ no walk outruns the graph distance, so cells
  farther than $i$ from the source are unoccupiable.
\item \emph{Backward} ($d(v,t)\le L-i$): the future constrains the
  present — from $v$ even the most direct route to $t$ costs $d(v,t)$
  moves, so a cell with $d(v,t)>L-i$ admits \emph{no} completion by the
  deadline, whatever the walk does afterwards. (At $i=L$ this test alone
  yields $D_L=\{t\}$; one step earlier only $t$ and its neighbours
  survive.)
\item \emph{Parity} ($i\equiv d(s,v) \bmod 2$): within reach is not
  enough — every move flips the checkerboard colour, so the step count
  and the distance must also agree modulo $2$.
\end{itemize}
The forward and backward tests are the two cones of a discrete
space--time ``light cone'' — reachable-from-$s$ and
can-still-reach-$t$ — and $D_i$ is their intersection restricted to the
correct colour class. Domains are accordingly narrow near both ends of
the horizon and widest in the middle. Deleting a cell that fails any
test removes no walk, since every walk satisfies all three by the
prefix/suffix argument of Proposition~\ref{prop:trim-exact}; the
backward family is also precisely what bucket elimination would
otherwise have to derive as resolvents (see below), while the forward
and parity families are stronger facts that the target-first elimination
would never produce.

\begin{proposition}[Exactness of the trimming]
\label{prop:trim-exact}
An assignment $(x_0,\dots,x_L)\in D_0\times\cdots\times D_L$ satisfies
$\Phi$ if and only if it is an absorbing $s$--$t$ walk: a four-neighbour
walk that starts at $s$, first reaches $t$ at some time
$\ell \le L$, and repeats $t$ thereafter. Moreover the trimming deletes
no such walk: the models of the trimmed formula and of the untrimmed
formula (with absorbing transitions and pinned endpoints) coincide.
Consequently \emph{every} simple $s$--$t$ path — of every length, the
maximum being $L=n-1$ — is a model, obtained by padding it with
repetitions of $t$; the remaining models are exactly the walks that
revisit a cell before reaching $t$.
\end{proposition}

\begin{proof}[Proof sketch]
($\Leftarrow$) Let $(x_i)$ be an absorbing walk with arrival time
$\ell$. For $i<\ell$ the prefix $x_0\cdots x_i$ is a genuine walk of
length $i$ from $s$, so $d(s,x_i)\le i$ and, by bipartiteness,
$i\equiv d(s,x_i) \pmod 2$; the suffix $x_i\cdots x_\ell$ reaches $t$ in
$\ell-i\le L-i$ moves, so $d(x_i,t)\le L-i$. For $i\ge\ell$, $x_i=t$ and
$i \ge \ell \ge d(s,t)$. Hence every position lies in its trimmed
domain, and consecutive positions satisfy $T$ by definition of a walk.
($\Rightarrow$) Conversely, a satisfying assignment moves along edges
until its first visit to $t$ (self-repetition is licensed only at $t$),
and $x_L=t$ because $D_L=\{t\}$; it is therefore an absorbing $s$--$t$
walk. Both directions never invoke the trimming tests negatively, which
proves the coincidence of model sets. Finally, a simple path padded at
$t$ is an absorbing walk, and an absorbing walk with pairwise-distinct
non-target positions is a padded simple path.
\end{proof}

\paragraph{Why this compiles with zero resolvents.}
The primal graph of $\Phi$ is the chain
$X_0 - X_1 - \cdots - X_L$ (induced width $1$), so bucket elimination
never couples more than two step variables. Its only derivations would
be the backward-reachability constraints
``$X_i$ must lie within distance $L-i$ of $t$'' propagated from the
pinned endpoint — precisely the \emph{backward} trimming test, which the
domains already encode; the forward and parity tests are stronger facts
that resolution would not derive at all. The trimmed input is therefore
its own compilation fixpoint. Measured compilation costs (single
worker, AMD EPYC 9335):

\begin{center}
\begin{tabular}{llrrrr}
\hline
grid & variant & clauses & resolvents & time & artifact \\
\hline
$12\times 12$ & untrimmed (waiting) & $20{,}448$ & $1{,}441$ & $38.2$\,s & $11$\,MB \\
$12\times 12$ & \textbf{trimmed}    & $8{,}839$  & $\mathbf{0}$ & $\mathbf{2.0}$\,\textbf{s} & $1.9$\,MB \\
$30\times 30$ & untrimmed (waiting) & $808{,}202$ & --- & unfinished after $19.8$\,h$^{*}$ & ($2.7$\,GB input) \\
$30\times 30$ & \textbf{trimmed}    & $379{,}315$ & $\mathbf{0}$ & $\mathbf{968}$\,\textbf{s} & $612$\,MB \\
\hline
\end{tabular}
\end{center}

\noindent
$^{*}$Snapshot at the time of writing: 37 of 900 buckets eliminated in
$19.8$\,h. The untrimmed $30\times 30$ compile proceeds at $\approx 30$
minutes per bucket through the $2(g{-}1)=58$-bucket backward-reachability
zone — deriving, one resolvent family at a time in width-$900$ literal
arithmetic, exactly the information the backward trimming test
precomputes — with a multi-day projected total. The trimmed encoding
eliminates this zone entirely at every grid size (zero resolvents at
$12\times 12$ and $30\times 30$ alike), leaving compilation dominated by
parsing and subsumption bookkeeping.

\paragraph{Worked example: $2\times 2$.}
Cells row-major: $s=0=(0,0)$, $1=(0,1)$, $2=(1,0)$, $t=3=(1,1)$;
$d(s,t)=2$, $L=3$. The trimmed domains are
\[
  D_0=\{0\},\qquad
  D_1=\{1,2\},\qquad
  D_2=\{3\},\qquad
  D_3=\{3\}.
\]
(E.g.\ $0\notin D_1$ by parity; $0\notin D_2$ because
$d(0,t)=2>L-2=1$; $t\notin D_1$ because $1<d(s,t)$.)
The formula has exactly two models,
\[
  (0,1,3,3) \qquad\text{and}\qquad (0,2,3,3),
\]
i.e.\ precisely the two simple $s$--$t$ paths — at this size the walk
relaxation is already exact, and the model count equals the path count.

\paragraph{Worked example: $3\times 3$.}
Cells $0..8$ row-major, $s=0$, $t=8$, $d(s,t)=4$, $L=8$. Distances from
$s$ are $(0,1,2,1,2,3,2,3,4)$ and to $t$ are $(4,3,2,3,2,1,2,1,0)$.
The trimmed domains are
\[
\begin{array}{ll}
  D_0=\{0\} &
  D_1=\{1,3\}\\
  D_2=\{0,2,4,6\} &
  D_3=\{1,3,5,7\}\\
  D_4=\{0,2,4,6,8\} &
  D_5=\{1,3,5,7,8\}\\
  D_6=\{2,4,6,8\} &
  D_7=\{5,7,8\}\\
  D_8=\{8\}
\end{array}
\]
— $29$ values in total against $81$ untrimmed (the parity test alone
halves each interior domain; the backward test additionally removes
$0$ from $D_6$ and $1,3$ from $D_7$; $t$ enters only from step
$d(s,t)=4$ on). Sample models:
\begin{itemize}
\item a geodesic simple path, padded at the target:
      $(0,1,2,5,8,8,8,8,8)$;
\item a maximal simple path using the full horizon (no padding):
      $(0,1,2,5,4,3,6,7,8)$;
\item a cyclic walk retained by \texttt{allow\_cycles} — it revisits
      $0$ and $3$ before arriving:
      $(0,1,4,3,0,3,6,7,8)$; loop erasure maps it to $(0,3,6,7,8)$.
\end{itemize}
All three satisfy every domain membership and transition. Excluded by
construction are: assignments that leave the target after reaching it
(absorption), assignments placing any cell off its parity class (e.g.\
$x_2=1$), jumps between non-adjacent cells, and prefixes that stray too
far to return (e.g.\ $x_6=0$: $d(0,t)=4>L-6$).

% ---------------------------------------------------------------------------
% Summary sentence for the experiments section, if wanted:
% The trimmed chain encoding represents every simple s-t path exactly,
% adds only loop-erasable cyclic walks, and is its own compilation
% fixpoint: COMPILE finds zero resolvents at every grid size tested,
% turning a multi-day 30x30 compilation into a 16-minute one.
% ---------------------------------------------------------------------------
